\clearpage
\section{Conclusion and Future Work}\label{section::Conclusion and Future Work}

This chapter concludes our work and outlines future research directions.

\subsection{Conclusion}
Addressing the pressing need for accurate crowd counting on highly dynamic railway platforms, we designed, implemented, and validated an end-to-end, high-performance multi-object tracking and counting system. The system employs a two-stage transfer-learned YOLOv11m as a high-precision head detector, a lightweight EfficientNet-B0 trained on scene-specific data as a robust appearance encoder, and integrates these into an advanced DeepSORT framework.
Our core contribution is the proposal and validation of a novel 6D physics-constrained Kalman motion model (Phys-6D). By injecting camera perspective geometry as a strong prior and introducing physically plausible motion constraints, Phys-6D significantly outperforms purely kinematic models on the key application metric. Compared to the standard 8D constant-velocity model (CV-8D) and the more complex 12D constant-acceleration model (CA-12D), Phys-6D maintains high tracking accuracy while substantially improving identity stability, achieving a mean absolute percentage counting error of just 2.97\% on our MOT-RPCH dataset.

To support this research, we constructed a comprehensive MOT-RailwayPlatformCrowdHead (MOT-RPCH) dataset containing 27 video sequences with 24,788 frames, 89,087 bounding boxes, and 885 unique human head identities. For MOT evaluation, we selected a representative subset of 20 video sequences with 18,548 frames, 73,799 bounding boxes, and 647 identities (\textbf{Tab.}~\ref{tab:20MOTValVideos}). This dataset represents one of the first large-scale, manually annotated multi-object tracking datasets specifically designed for railway platform crowd scenarios, providing a valuable benchmark for future research in this domain.
In summary, beyond delivering a practical solution ready for intelligent transportation scenarios, this work provides compelling evidence that, in vision tasks governed by clear physical laws, fusing domain knowledge with data-driven learning is a decisive path toward higher performance and stronger robustness.

\subsection{Future Work}

Building on these results and the identified limitations, we highlight the following forward-looking directions to further enhance intelligence, generalization, and applicability.

\subsubsection{Model Adaptivity and Generalization}
To reduce reliance on pre-calibrated intrinsics, a first priority is online camera self-calibration—estimating intrinsics dynamically from trajectories and size changes in video alone. Achieving this would greatly improve practicality and enable true plug-and-play deployment. In parallel, we plan to expand the dataset to include diverse weather, lighting, and platform layouts, and investigate domain adaptation/generalization to train a single, universal model that transfers to unseen railway environments.

\subsubsection{Toward End-to-End and Joint Learning Frameworks}
Our current pipeline is modular (detect–track). While interpretable and debuggable, independently optimized modules may be suboptimal globally. Recent end-to-end transformer-based trackers (e.g., MOTR, TrackFormer) show promise. We will explore embedding physics constraints into such frameworks—for example, physics-aware attention or queries that implicitly learn and exploit scene geometry—enabling joint optimization of detection, tracking, and physical modeling within a unified network.

\subsubsection{Multi-Modal Fusion}
Pure RGB systems degrade at night or in adverse weather. Modern hubs deploy heterogeneous sensors. We will study fusion with LiDAR and thermal cameras. LiDAR offers illumination-invariant 3D depth; thermal enables perception in darkness. With principled fusion algorithms, we aim to build an all-weather, round-the-clock tracker with exceptional robustness.

\subsubsection{From Counting to Business Intelligence and Resource Optimization}
Accurate counting underpins higher-level analytics. Our long-term vision is to move beyond “counting heads” toward a data-driven operational paradigm. The system’s high-quality spatiotemporal density data constitute a valuable intelligence asset, revealing population distributions, flows, and dwell hotspots across lines and platforms, thereby informing resource allocation.
Commercial layout and advertising: mining long-horizon traffic uncovers “golden spots” with maximal density and dwell time, providing quantitative evidence for retail placement and dynamic, targeted digital advertising with measurable impact.
Planning and strategic decisions: the accumulated data serve as the foundation for long-term planning (e.g., line-specific patterns guide investment; sustained monitoring of space utilization informs renovations), ensuring each investment advances both commercial value and service quality.







